\section{Ostrowski data and the complete repetition table}
\label{sec:repetition}

Let $0<\alpha<1$ be irrational, with continued fraction
$\alpha=[0;a_1,a_2,\ldots]$.  Its continuant denominators are normalized by
\begin{equation}
  q_{-1}=0,
  \qquad q_0=1,
  \qquad q_{n+1}=a_{n+1}q_n+q_{n-1}.
  \label{eq:continuants}
\end{equation}
A digit sequence $(b_i)_{i\geq1}$ is \emph{legal} if
\begin{equation}
  0\leq b_1\leq a_1-1,
  \qquad
  0\leq b_i\leq a_i\quad(i\geq2),
  \qquad
  b_{i+1}=a_{i+1}\Longrightarrow b_i=0.
  \label{eq:legality}
\end{equation}
Following the formal-intercept convention of \cite[Section~2]{Wojcik2025v1},
put
\begin{equation}
  \rho_n=\sum_{i=0}^{n-1}b_{i+1}q_i,
  \qquad
  \rho_{n+1}=\rho_n+b_{n+1}q_n.
  \label{eq:rho-recurrence}
\end{equation}
Legality gives $0\leq\rho_n<q_n$.  The infinite digit data specify a formal
intercept $\rho$ and hence a Sturmian word $x=T^\rho(c_\alpha)$; no numerical
convergence of the displayed formal sum is being asserted.

For an infinite word $x$, let $r(x,m)$ be the largest $k$ for which the first
$k$ length-$m$ factors
\[
  \mathbb P_m(x),\mathbb P_m(Tx),\ldots,
  \mathbb P_m(T^{k-1}x)
\]
are pairwise distinct.  This is the convention used in
\cite[Definition~1.1.2]{Wojcik2025v1}; it differs by an additive $m$ from the
prefix-length convention in \cite{BugeaudKim2019}.  For a Sturmian word,
$r(x,m)\leq m+1$, and the standard repetition identity becomes
\begin{equation}
  \dio(x)=1+\limsup_{m\to\infty}\frac{m}{r(x,m)}.
  \label{eq:dio-repetition}
\end{equation}

We now record the finite interface used in the proof.  Write
\begin{equation}
\begin{gathered}
 q=q_n,\quad p=q_{n-1},\quad Q=q_{n+1},\\
 s=\rho_{n-1},\quad t=\rho_n,\quad u=\rho_{n+1},
 \qquad b=b_{n+1},\quad a=a_{n+1}.
\end{gathered}
\label{eq:table-shorthand}
\end{equation}
The integer interval belonging to level $n$ is
$I_n=[q-1,Q-2]$.  Each row below has the form
``$[A,B];V$'', meaning $r(x,m)=V$ for every integer $m\in[A,B]$.
A row is discarded exactly when $A>B$.  Its last column records the ratio
obtained from the right endpoint after the finite constants disappear in a
fixed-phase limit.

\begingroup
\scriptsize
\setlength{\tabcolsep}{3.2pt}
\renewcommand{\arraystretch}{1.16}
\begin{longtable}{@{}>{\centering\arraybackslash}p{0.055\textwidth}
  >{\RaggedRight\arraybackslash}p{0.255\textwidth}
  >{\RaggedRight\arraybackslash}p{0.415\textwidth}
  >{\RaggedRight\arraybackslash}p{0.19\textwidth}@{}}
\caption{The complete plateau interface.  Case~7 contains the correction
$t=\rho_n$ in place of the printed $s=\rho_{n-1}$.
\label{tab:eight-cases}}\\
\toprule
Case & Hypotheses & Candidate plateau $[A,B]$; value $V$ & Phase ratio $B/V$ \\
\midrule
\endfirsthead
\multicolumn{4}{c}{\tablename\ \thetable\ (continued)}\\
\toprule
Case & Hypotheses & Candidate plateau $[A,B]$; value $V$ & Phase ratio $B/V$ \\
\midrule
\endhead
\midrule
\multicolumn{4}{r}{Continued on the next page}\\
\endfoot
\bottomrule
\endlastfoot

1 & $b=0$, $b_{n+2}=a_{n+2}$
& $[q-1,Q-2];\ q$
& $Q/q$ \\
\addlinespace

2 & $b=0$, $b_{n+2}\ne a_{n+2}$, $b_n=0$
& $[q-1,Q-s-2];\ q$\newline
  $[Q-s-1,Q-2];\ Q-s$
& $(Q-s)/q$\newline
  $Q/(Q-s)$ \\
\addlinespace

3 & $b=0$, $b_{n+2}\ne a_{n+2}$, $a\ne1$
& $[q-1,Q-t-2];\ q$\newline
  $[Q-t-1,Q-2];\ Q-t$
& $(Q-t)/q$\newline
  $Q/(Q-t)$ \\
\addlinespace

4 & $b=0$, $b_{n+2}\ne a_{n+2}$, $a=1$, \mbox{$b_n\ne0$}
& $[q-1,Q-2];\ Q-t=q+p-t$
& $Q/(Q-t)$ \\
\addlinespace

5 & $0<b<a-1$
& $[q-1,Q-u-2];\ q$\newline
  $[Q-u-1,Q-bq-2];\ Q-u$\newline
  $[Q-bq-1,Q+q-u-2];\ Q-bq$\newline
  $[Q+q-u-1,Q-2];\ Q+q-u$
& $(Q-u)/q$\newline
  $(Q-bq)/(Q-u)$\newline
  $(Q+q-u)/(Q-bq)$\newline
  $Q/(Q+q-u)$ \\
\addlinespace

6 & $b=a-1$, $a\ne1$, $b_n=0$
& $[q-1,q+p-s-2];\ q$\newline
  $[q+p-s-1,q+p-2];\ q+p-s$\newline
  $[q+p-1,2q+p-s-2];\ q+p$\newline
  $[2q+p-s-1,Q-2];\ 2q+p-s$
& $(q+p-s)/q$\newline
  $(q+p)/(q+p-s)$\newline
  $(2q+p-s)/(q+p)$\newline
  $Q/(2q+p-s)$ \\
\addlinespace

\textbf{7} & \mbox{$b=a-1$, $a\ne1$, $b_n\ne0$};\newline \textbf{corrected}
& $[q-1,q+p-2];\ q+p-t$\newline
  $[q+p-1,2q+p-t-2];\ q+p$\newline
  $[2q+p-t-1,Q-2];\ 2q+p-t$
& $(q+p)/(q+p-t)$\newline
  $(2q+p-t)/(q+p)$\newline
  $Q/(2q+p-t)$ \\
\addlinespace

8 & $b=a$ (hence $b_n=0$ and $t=s$)
& $[q-1,q+p-s-2];\ p$\newline
  $[q+p-s-1,Q-2];\ q+p-s$
& $(q+p-s)/p$\newline
  $Q/(q+p-s)$ \\
\end{longtable}

\endgroup

The table is the content of the complete case split in
\cite[Theorem~2.2.1]{Wojcik2025v1}, transcribed with finite endpoint constants
intact, except for the explicitly marked correction in case~7.  Cases~2 and~3
can overlap; then $b_n=0$ gives $t=s$, so their rows agree.  In case~8,
legality likewise forces $b_n=0$ and hence $t=s$.  Thus the apparent overlap
does not increase the number of distinct plateau rows.

\begin{proposition}[Finite-table interface]
\label{prop:finite-interface}
For all sufficiently large $n$, the applicable rows in
\cref{tab:eight-cases} cover $I_n$, and $r(x,m)$ has the displayed constant
value on each nonempty row.  Consequently, if a row is $[A_n,B_n];V_n$, its
contribution to the limsup in \eqref{eq:dio-repetition} is determined by
$B_n/V_n$.
\end{proposition}

\begin{proof}
Coverage and the plateau values, apart from the case-7 subscript, are precisely
the cited source theorem; the missing correction is proved independently in
\cref{lem:case-seven}.  Since $V_n$ is a positive integer and is constant on a
nonempty row, $m/V_n$ increases with $m$.  Its maximum on that row is therefore
$B_n/V_n$.
\end{proof}
